\begin{pastexam}[label={px:nchu-114-dsa-14}]{中興資工所甲組 114 學年度 單選題第 14 題}
（單選；每題 3 分，未答題得 0 分，答錯不倒扣）

Which of the following data structures is best suited for implementing the undo operation in a text editor?

(A)~Stack \qquad (B)~Queue \qquad (C)~Binary Search Tree \qquad (D)~Hash Table
\end{pastexam}

\begin{pastsoln}
答案是 (A)。

復原（undo）要先撤銷「最近一次」的操作，再撤銷前一次，順序是後做的先撤銷，正是堆疊的後進先出
（資結書第~6.1~節）。 每做一個操作就 push 一筆紀錄，按復原就 pop 出最上面那筆並反向執行，兩者都是 $O(1)$。
（重做 redo 通常再用第二個堆疊，接住被復原的紀錄。）
\begin{itemize}
\item (B) Queue：先進先出，會先撤銷最早的操作，順序相反。 錯。
\item (C) BST：依鍵值排序，和操作的時間順序無關。 錯。
\item (D) Hash Table：用鍵值快速查詢，沒有保存先後順序。 錯。
\end{itemize}
\end{pastsoln}
